\section{Lifting the product to the complex cycle algebra}\label{sec:product}

\subsection{Stabilization of the integers}
At cutoff $\Lambda$, let $N^\Lambda_{\gamma,h}$ be the sum of the signs
of the elementary factors of genus $h$ and core class $\gamma$ in
Proposition~\ref{prop:finite}. These integers are zero outside a finite
set of pairs. Suppose two decompositions, possibly one obtained by
truncating a decomposition at a higher bound, give different integers.
Cancel their common lower-degree factors and take a class $\beta$ in
the first degree of disagreement. Because every nonconstant term of a
factor in class $\gamma$ has degree at least that of $\gamma$, comparison
of the coefficients of $Q^\beta$ gives
\[
           \sum_h (N^\Lambda_{\beta,h}-\widetilde N_{\beta,h})\,
                       \Phi(x_{h,1})=0 .
\]
Here we used $a_{h,1}=x_{h,1}$. Theorem~\ref{thm:independence} gives
equality in every genus, a contradiction. The same induction proves
that $N^\Lambda_{\beta,h}=0$ when $\beta$ is ineffective for the
original projective structure. The compact series has coefficient
zero there, and sums of effective lower labels are effective.

For the chosen polarization denote the stabilized counts by
$n_{\beta,h}$. Taking $\Lambda$ equal to the degree of $\beta$ proves
finite genus support at $\beta$. Thus Proposition~\ref{prop:finite}
has already given the product of Theorem~\ref{thm:main} after
application of $\Phi$. We need only multiplicity-one independence to
lift this identity.

\subsection{Primitive classes}
Let $P=\QQ[X_{h,m}:h\ge0,\ m\ge1]$ and $X_{h,0}=1$. The map
$p:P\to H$ taking $X_{h,m}$ to $x_{h,m}$ is surjective by the
geometric generation theorem \cite[Theorem~6.13]{Pardon}.
The Chern numbers of the generators of virtual dimension zero are
nonnegative, so only generators of Chern number zero can enter
this summand. The scalar classes respect the coproduct $\Delta$ from
addition of cycles and pullback $\rho_d$ under multiplication by $d$.
Their formulas are
\[
 \Delta X_{h,m}=\sum_{r+s=m}X_{h,r}\otimes X_{h,s},
 \qquad
 \rho_d X_{h,m}=
 \begin{cases}X_{h,m/d}& d\mid m,\\ 0& d\nmid m.\end{cases}
\]
For the scalar classes, apply the division formula and the diagonal
calculation of \cite[\S4, in particular Lemma~4.13]{Pardon}. Each term
with a nonzero shift has one tensor factor of negative virtual dimension
and thus vanishes by \cite[Theorem~6.13]{Pardon}. The same vanishing
removes the shifted terms in the balanced calculation
\cite[\S3]{PaperIII}. Thus both operation formulas use only the
geometric generation and dimension statements on the complex side. The augmentation on both algebras is
restriction to the empty cycle.

In the triangular coordinates
$v_{h,m}=[q^m]\log(\sum_{j\ge0}X_{h,j}q^j)$ the generators are
primitive and $\rho_d v_{h,m}=v_{h,m/d}$ for $d\mid m$, zero otherwise.
Suppose the kernel of $p$ contains a nonzero polynomial $f$ of smallest
total polynomial degree in these coordinates. Its constant term is zero
by augmentation. Its reduced coproduct lies in the tensor square of
polynomials of strictly lower degree. The map $p$ is injective on
this vector space and remains so on its tensor square. Since the
reduced coproduct maps to zero, $f$ is primitive. Over $\QQ$ an
additive polynomial is linear, by differentiating $f(v+w)=f(v)+f(w)$
at $w=0$. Applying $\rho_j$ for the largest multiplicity $j$ occurring
in $f$ gives a nonzero formal combination of the $X_{h,1}$ in the
kernel. Their images even under $\Phi p$ are independent by
Theorem~\ref{thm:independence}. Hence $p$ is an isomorphism.

\begin{lemma}\label{lem:primitive-division}
Let $\operatorname{Prim}(H)$ denote the primitive subspace for $\Delta$.
Then
\[
  \operatorname{Prim}(H)\cap\bigcap_{d>1}\ker\rho_d
      =\bigoplus_{h\ge0}\QQ x_{h,1}.
\]
\end{lemma}
\begin{proof}
In the coordinates $p(v_{h,m})$ a primitive is an additive polynomial,
hence linear. If its largest multiplicity is $j>1$, applying $\rho_j$
leaves exactly the coefficients at multiplicity $j$, now on the
linearly independent classes $x_{h,1}$. Conversely, each $x_{h,1}$
is primitive and is killed by every $\rho_d$ with $d>1$.
\end{proof}

\subsection{The first coefficient of disagreement}
Write $F$ for either the compact series or the product with the
stabilized exponents, formed in $H\llbracket\Gamma_X\rrbracket$.
Each has $F_0=1$ and satisfies
\begin{equation}\label{eq:ops}
 \Delta F_\beta=\sum_{\gamma+\delta=\beta}F_\gamma\otimes F_\delta,
 \qquad
 \rho_d F_\beta=\sum_{\substack{\gamma\in\Gamma_X\\d\gamma=\beta}}F_\gamma .
\end{equation}
For the compact series take inverse images of the entire class locus
under addition and multiplication. For the other series use the
formulas on the balanced elements, multiply, and take inverses for
negative exponents. The substitution $Q^\gamma\mapsto Q^{d\gamma}$
is multiplicative even when several exponents have the same image.
Each effective root occurs separately in \eqref{eq:ops}.

If $R_\beta$ is a nonzero difference at a least ample degree, the
terms with two positive summands cancel in its coproduct. It is
primitive. For $d>1$ all the roots on the right of \eqref{eq:ops}
have smaller degree, so $\rho_d R_\beta=0$. By
Lemma~\ref{lem:primitive-division}, $R_\beta$ is a linear combination
of the $x_{h,1}$.
The finite-energy identity says $\Phi(R_\beta)=0$.
Theorem~\ref{thm:independence} now forces $R_\beta=0$.
This proves the product formula. Uniqueness follows by comparing the
first degree of disagreement between any two systems of exponents.
A second polarization yields an identity in the same coefficients. For either polarization there are finitely many effective sublabels
of any given $\beta$, and their possible decompositions are finite.
Uniqueness therefore identifies the two systems.

\subsection{Identification with the elementary-cluster coefficients}
It remains to identify the integers already obtained with those of
the connected Gromov--Witten expansion. Apply fixed-target connected
counts to the \emph{geometric decompositions} of
\S\ref{sec:clusters}. For a compact open insertion $u$ let
$\epsilon(u)$ be its value at the empty cycle, and let $D_k(u)$ be
its localized positive-energy connected genus-$k$ count. The same
open part of the stable-map space computes this count on excision
into a target neighborhood. For the compact symplectic targets
and paths in that section the localized counts are additive on
partitions and invariant under transport
\cite[Lemma~3.8 and Proposition~3.10]{DIW}. On disjoint targets a
connected map lies in just one of the two, hence
\[
                D_k(uv)=D_k(u)\epsilon(v)+\epsilon(u)D_k(v).
\]
Extend the equality coefficientwise with the integral class labels.
Each of our normalized series has empty-cycle value one as a series.
Thus connected counts add under their multiplication and change sign
on inversion, which can equally be calculated in the truncated
series algebra. This applies along the exact fold, subdivision and
endpoint relations used to obtain each finite product.

For genus $h$ the elementary connected contribution is
\[
 G_h(q,t)=\log\!\left(1+
 \sum_{m\ge1}q^m\sum_{\lambda\vdash m}
       \prod_{\square\in\lambda}
           \bigl(2\sin(\ell(\square)t/2)\bigr)^{2h-2}\right),
\]
where $\ell(\square)$ is the hook length \cite[Proposition~3.3]{IP}.
Thus the signed occurrences of each factor give the coefficients of
$G_h(Q^\beta,t)$ in the connected expansion of $X$. Uniqueness is
\cite[Theorem~3.50]{DIW}. At the first degree of disagreement, the
least remaining genus is distinguished by
$[q]G_h(q,t)=t^{2h-2}+O(t^{2h})$. This identifies the integers of
\eqref{eq:product} and completes the proof of Theorem~\ref{thm:main}.
